% Chapter 2, Topic _Linear Algebra_ Jim Hefferon
%  http://joshua.smcvt.edu/linearalgebra
%  2001-Jun-11
\topic{晶体}
\index{crystals|(}
人人都注意到，
食盐\index{salt}
是小立方体形状的。
\begin{center}
  \includegraphics[height=1.25in]{vs/pix/salt.jpg} %1.25in tall
\end{center}
这种整齐的外观源于整齐的内部\Dash 
原子的排布方式也是立方形的，
这些立方体整齐地按行与列堆叠，而食盐的晶面
往往恰好就是立方体的最外层。
下面展示了一个由原子构成的立方体。
食盐是氯化钠，图中
小球表示钠原子，大球表示氯原子。
为了简化视图，图中只画出了前侧、
顶部与右侧的钠原子和氯原子。
\begin{center}
  \includegraphics{vs/mp/ch2.1}
\end{center}
我们在上图中看到的盐粒
是这种基本单元的多次重复。
像食盐这样的固体，
若具有规则的内部结构，就称为\definend{晶体}。

我们可以把注意力限制在前表面上。
在那里，一个正方形被重复了很多次，形成原子的晶格。
\begin{center}
  \includegraphics{vs/mp/ch2.9}
\end{center}
每个正方形晶胞沿边的距离约为
$3.34$~埃
（一埃是 $10^{-10}$~米）。
当我们要指称晶格中的原子时，这个数用起来很不方便，
于是我们
取正方形的边长作为单位。
也就是说，我们自然地采用下面这个基。
\begin{equation*}
  \sequence{\colvec{ 3.34 \\ 0},\colvec{0 \\ 3.34}}
\tag*{}\end{equation*}
现在，比如说，我们可以把上面晶格图中右上角的原子
描述为
$3\vec{\beta}_1+2\vec{\beta}_2$，而不是“向右 $10.02$~埃、
向上 $6.68$~埃”。

来自日常经验的另一种晶体是铅笔芯。
它是\definend{石墨}，\index{graphite} 
由碳原子按下图的形状排列而成。
\begin{center}  %graphite
  \includegraphics{vs/mp/ch2.10}
\end{center}
这是石墨的一个单独平面，称为\definend{石墨烯}。
一块石墨由许多这样的平面层叠而成。
平面之间的化学键
比平面内部的键弱得多，这解释了铅笔为什么能写字\Dash 
石墨可以被剪切，使这些平面滑脱
并留在纸上。

我们可以用如下方法得到一个方便的长度单位：把六边形环分解成三个区域，
它们都是这个\definend{原胞} 的旋转。\index{crystals!unit cell} 
\begin{center}  %graphite
  \includegraphics{vs/mp/ch2.11}
\qquad\qquad
\includegraphics{vs/mp/ch2.12}
\end{center}
构成该原胞各条边的向量
组成了一个方便的基。
底边与斜边的长度都是 $1.42$~埃，
所以
\begin{equation*}
  \sequence{\colvec{1.42 \\ 0}, \colvec{0.71 \\ 1.23}}
\tag*{}\end{equation*}   
是一个很好的基。

另一种由碳形成的、我们熟悉的晶体是金刚石。\index{diamond}
与食盐一样，它也由立方体构成，但每个立方体
内部的结构更加复杂。
除了每个角上各有一个碳原子之外，
\begin{center}
  \includegraphics{vs/mp/ch2.13}
\end{center}
每个面的中央还有碳原子。
\begin{center}
  \includegraphics{vs/mp/ch2.14}
\end{center}
（为了清楚地显示新出现的面心碳原子，
角上的碳原子被缩小成了点。）
立方体内部还有四个碳原子，
其中两个位于自底面向上四分之一处，
另两个位于自顶面向下四分之一处。
\begin{center}
  \includegraphics{vs/mp/ch2.15}  
\end{center}
（与前面一样，之前出现过的碳原子在这里被缩小成了点。）
立方体任意一条棱的长度都是 $2.18$~埃。
于是，描述碳原子的位置
以及它们之间化学键的一个自然基就是下面这个。
\begin{equation*}
  \sequence{\colvec{2.18 \\ 0 \\ 0}, 
            \colvec{0 \\ 2.18 \\ 0}, 
            \colvec{0 \\ 0 \\ 2.18}}
\tag*{}\end{equation*}   

这里的例子表明，
晶体结构相当复杂，需要某种有条理的体系来给出原子的位置
以及它们如何化学键合。
实现这种组织的一个工具就是方便的基。
基的这个应用虽然简单，却展示了一个
科学背景，
这个想法在其中自然地出现。
% The work in this chapter just takes this simple idea and develops it.

\begin{exercises}
  \item 
    一粒盐的一个面上有多少个基本区域？
   （用直尺我们可以估计出，该面是一个边长为 $0.1$~cm 的正方形。）
   \begin{answer}
     每个基本单元是 $3.34\times 10^{-10}$~cm，所以这样的单元约有
     $0.1/(3.34\times 10^{-10})$ 个。
     这给出 $2.99\times 10^{8}$，因此差不多有
     $300,000,000$（三亿）个区域。     
   \end{answer}
  \item 
    在石墨的图中，假设我们关注这样一点：它位于从原点出发
    向右 $5.67$~埃、向上 $3.14$~埃处。
    \begin{exparts}
      \partsitem 用针对石墨给出的基来表示该点。
      \partsitem 这个点距离原点有多少个六边形？
      \partsitem 用第二个基来表示该点，
        其中第一个基向量与原来相同，但第二个基向量
        与第一个垂直（沿平面向上）且长度相同。
    \end{exparts} 
    \begin{answer}
      \begin{exparts}
        \partsitem 我们求解
          \begin{equation*}
            c_1\colvec{1.42 \\ 0}+c_2\colvec{0.71 \\ 1.23}
              =\colvec{5.67 \\ 3.14}
            \quad\implies\quad
            \begin{linsys}{2}
              1.42c_1  &+  &0.71c_2  &=  &5.67  \\
                       &   &1.23c_2  &=  &3.14
            \end{linsys}
          \tag*{}\end{equation*}
          解得 $c_1\approx 2.74$ 与 $c_2\approx 2.55$。
        \partsitem 下面是这个点在晶格中的定位。
          在左图中，原胞上叠加了
          两个基向量 $\vec{\beta}_1$ 与 $\vec{\beta}_2$， 
          以及一个标出偏移量
          $2.55\vec{\beta}_1+2.72\vec{\beta}_2$ 的方框。
          右图显示的是它在晶体晶格内部出现的位置，这里取
          左下角那个六边形的左下顶点作为原点。
          \begin{center}  %graphite
            \includegraphics{vs/mp/ch2.17}
            \qquad
            \includegraphics{vs/mp/ch2.18}
       \end{center}
       所以这个点向右移过两列六边形，
       向上移过一行六边形。
      \partsitem 这第二个基
          \begin{equation*}
            \sequence{\colvec{1.42 \\ 0},\colvec{0 \\ 1.42}}
          \tag*{}\end{equation*}
          使计算更容易
          \begin{equation*}
            c_1\colvec{1.42 \\ 0}+c_2\colvec{0 \\ 1.42}
              =\colvec{5.67 \\ 3.14}
            \quad\implies\quad
            \begin{linsys}{2}
              1.42c_1  &   &         &=  &5.67  \\
                       &   &1.42c_2  &=  &3.14
            \end{linsys}
          \tag*{}\end{equation*}
          （我们得到 $c_2\approx 2.21$ 与 $c_1\approx 3.99$），但它似乎
          与我们正在研究的物理结构没有多大关系。
      \end{exparts}
    \end{answer}
  \item 
    用基以及用埃两种方式给出金刚石立方体中各原子的位置。
    \begin{answer}
      用基表示时，角上原子的位置是
      $(0,0,0)$、$(1,0,0)$、\ldots{}、$(1,1,1)$。
      面上原子的位置是 $(0.5,0.5,1)$、$(1,0.5,0.5)$、
      $(0.5,1,0.5)$、$(0,0.5,0.5)$、$(0.5,0,0.5)$ 以及 $(0.5,0.5,0)$。
      自顶面向下四分之一处的原子的位置
      是 $(0.75,0.75,0.75)$ 与 $(0.25,0.25,0.25)$。
      自底面向上四分之一处的原子
      位于 $(0.75,0.25,0.25)$ 与 $(0.25,0.75,0.25)$。
      换算成埃很容易。
    \end{answer}
  \item 
    这说明了我们如何从一种物质结晶的形状
    计算出原胞的尺寸
    （\cite{Ebbing}，第~462~页）。
    \begin{exparts}
      \partsitem 回顾一下，每摩尔有 $6.022\times 10^{23}$ 个原子
        （这就是阿伏伽德罗数）。
        由此，再加上铂每摩尔质量为 $195.08$ 克
        这一事实，计算每个原子的质量。
      \partsitem 铂以面心立方晶格的形式结晶，
        每个格点上都有一个原子，
        也就是说，它看起来像上面金刚石晶体插图中的
        中间那幅。
        求每个原胞中铂原子的个数
        （提示：~把位于单个晶胞内部的各部分铂原子
        所占的份额加起来）。
      \partsitem 由此求出一个原胞的质量。
      \partsitem 铂晶体的
        密度为每立方厘米 $21.45$ 克。
        由此以及一个原胞的质量，计算
        一个原胞的体积。
      \partsitem 求每条棱的长度。
      \partsitem 描述一个自然的三维基。
    \end{exparts}
    \begin{answer}
      \begin{exparts}
        \partsitem $195.08/6.02\times 10^{23}=3.239\times 10^{-22}$
        \partsitem 每个面中央的铂原子
          由两个立方体分摊，所以到这里共有 $6/2=3$ 个原子。
          每个角上的原子由八个立方体分摊，所以
          又增加 $8/8=1$~个原子，因此总数是 $4$。
        \partsitem $4\cdot 3.239\times 10^{-22}=1.296\times 10^{-21}$
        \partsitem $1.296\times 10^{-21}/21.45=6.042\times 10^{-23}$ 立方
           厘米
        \partsitem $3.924\times 10^{-8}$ 厘米。
        \partsitem 
            $\sequence{\colvec{3.924\times 10^{-8} \\ 0 \\ 0},
                       \colvec{0 \\ 3.924\times 10^{-8} \\ 0},
                       \colvec{0 \\ 0 \\ 3.924\times 10^{-8}}}$
      \end{exparts}
    \end{answer}
\end{exercises}
\index{crystals|)}
\endinput
